% Vietnamese translation for OI Public Library
% Source-PDF: 比赛 CONTEST/2021“MINIEYE杯”中国大学生算法设计超级联赛/2021“MINIEYE杯”中国大学生算法设计超级联赛（5）/2021“MINIEYE杯”中国大学生算法设计超级联赛（5）-题目集.pdf
\documentclass[11pt]{article}
\usepackage{oipl-vn}

\title{Đề bài HDU Multi-University 2021, Contest 5}
\author{}
\date{}

\begin{document}
\maketitle

\origtitle{2021 HDU Multi-University Training Contest 5}
\sourcepath{contest/hdu-2021-minieye-05-problems.pdf}

\section{Miserable Faith}

\paragraph{Input file:} standard input
\paragraph{Output file:} standard output
\paragraph{Time limit:} 2 seconds
\paragraph{Memory limit:} 128 megabytes

Cho một cây gồm \(n\) đỉnh đánh số \(1,\ldots,n\), gốc là đỉnh \(1\). Ban
đầu màu của đỉnh thứ \(i\) là \(i\). Với mỗi cạnh nối hai đỉnh \(u,v\), trọng
số cạnh bằng \(1\) nếu màu của \(u\) khác màu của \(v\), và bằng \(0\) nếu
hai màu giống nhau.

Ký hiệu đường đi giữa \(u\) và \(v\) là \((u,v)\). Gọi \(d(u,v)\) là số cạnh
trên đường đi, và \(\operatorname{dist}(u,v)\) là tổng trọng số các cạnh trên
đường đi đó. Nếu \(u\) là tổ tiên của \(v\), và \(v\) là tổ tiên của \(w\),
thì \(u\) là tổ tiên của \(w\). Ta viết \(v\in\operatorname{subtree}(u)\) nếu
\(u\) là tổ tiên của \(v\).

Định nghĩa
\[
  \operatorname{FakeDeep}_u=
  \max\{d(i,u)\mid 1\le i\le n,\ i\text{ là tổ tiên của }u,
  \operatorname{dist}(i,u)=0\}.
\]
Có \(m\) truy vấn, thuộc một trong bốn loại:
\begin{itemize}
  \item \texttt{1 u c}: đổi màu mọi đỉnh trên đường \((u,1)\) thành \(c\);
        màu \(c\) khác mọi màu đang có trong cây.
  \item \texttt{2 u v}: in \(\operatorname{dist}(u,v)\).
  \item \texttt{3 u}: in
        \(\sum_{i\in\operatorname{subtree}(u)}\operatorname{dist}(i,u)\).
  \item \texttt{4}: in \(\sum_{i=1}^n\operatorname{FakeDeep}_i\).
\end{itemize}

\subsection*{Input}

Dữ liệu vào gồm \(T\) bộ dữ liệu (\(1\le T\le 10\)). Với mỗi bộ dữ liệu, dòng đầu
chứa hai số nguyên \(n,m\) (\(1\le n,m\le 10^5\)). \(n-1\) dòng tiếp theo mỗi
dòng chứa hai đỉnh \(u,v\), biểu thị một cạnh. Sau đó là \(m\) dòng truy vấn,
có dạng \texttt{1 u c}, \texttt{2 u v}, \texttt{3 u} hoặc \texttt{4}, với
\(1\le u,v\le n\), \(1\le c\le 10^6\).

\subsection*{Output}

Với mỗi truy vấn loại \(2,3,4\), in một dòng chứa đáp án.

\subsection*{Sample}

\begin{verbatim}
Input
1
7 7
1 2
1 6
6 7
3 4
1 3
3 5
1 4 8
1 7 4
4
2 5 7
3 3
1 2 3
4

Output
4
2
1
3
\end{verbatim}

\section{String Mod}

\paragraph{Time limit:} 5.5 seconds
\paragraph{Memory limit:} 256 megabytes

Xét mọi xâu độ dài \(L\) tạo từ \(k\) chữ cái thường đầu tiên
(\(2\le k\le 26\)); tổng số xâu là \(k^L\). Với mỗi cặp
\((i,j)\), \(0\le i,j\le n-1\), cần đếm số xâu có \(p\) ký tự \texttt{a} và
\(q\) ký tự \texttt{b}, thỏa
\[
  p\equiv i\pmod n,\qquad q\equiv j\pmod n.
\]
Hãy in ma trận \(A\), trong đó \(A[i][j]\) là đáp án cho cặp \((i,j)\).

\subsection*{Input}

Dữ liệu vào gồm \(T\le 10\) bộ dữ liệu. Mỗi bộ dữ liệu gồm ba số
\[
  k,\ L,\ n
\]
với \(2\le k\le 26\), \(1\le L\le 10^{18}\), \(0\le n\le 500\).

\subsection*{Output}

Với mỗi bộ dữ liệu, in \(n\) dòng, mỗi dòng \(n\) số. Đáp án lấy modulo
\[
  P=10^9+9,
\]
với điều kiện \(n\mid(P-1)\). Tổng các \(n\) trong input không vượt quá
\(2000\).

\subsection*{Sample}

\begin{verbatim}
Input
2 1 3

Output
0 1 0
1 0 0
0 0 0
\end{verbatim}

\section{VC Is All You Need}

\paragraph{Time limit:} 1 second
\paragraph{Memory limit:} 256 megabytes

Trong mặt phẳng hai chiều, với ba điểm bất kỳ, ta có thể vẽ một đường thẳng
tách các điểm theo màu sao cho các điểm cùng màu nằm về cùng một phía, với
bất kỳ cách tô màu xanh/đỏ nào cho ba điểm. Với bốn điểm thì tính chất này
không còn đúng.

Tổng quát, trong không gian Euclid thực \(k\) chiều \(\mathbb{R}^k\), hãy hỏi
có tồn tại \(n\) điểm sao cho với mỗi trong \(2^n\) cách tô màu, luôn tồn tại
một siêu phẳng \(k-1\) chiều tách hai màu hay không.

\subsection*{Input}

Dữ liệu vào gồm \(T\) bộ dữ liệu (\(1\le T\le 10^5\)). Mỗi bộ dữ liệu gồm hai số
\[
  n,k\in[2,10^{18}].
\]

\subsection*{Output}

In \texttt{Yes} nếu có thể tìm được tập điểm như vậy, ngược lại in \texttt{No}.

\subsection*{Sample}

\begin{verbatim}
Input
3
2 2
3 2
4 2

Output
Yes
Yes
No
\end{verbatim}

\section{Another String}

\paragraph{Time limit:} 1 second
\paragraph{Memory limit:} 256 megabytes

Khoảng cách giữa hai xâu cùng độ dài là số vị trí mà hai ký tự khác nhau. Nếu
khoảng cách giữa hai xâu không vượt quá \(k\), ta nói hai xâu thỏa
\(k\)-matching.

Cho xâu \(S\) độ dài \(n\) và số nguyên \(k\). Với mỗi \(i\in[1,n-1]\), tách
\[
  A=S[1,i],\qquad B=S[i+1,n].
\]
Trong tất cả cặp xâu tạo bởi một xâu con không rỗng của \(A\) và một xâu con
không rỗng của \(B\), đếm số cặp thỏa \(k\)-matching.

\subsection*{Input}

Dữ liệu vào gồm \(T\le 60\) bộ dữ liệu. Mỗi bộ dữ liệu gồm \(n,k\)
\[
  2\le n\le 3000,\qquad 0\le k\le 3000,
\]
sau đó là xâu \(S\). Xâu chỉ gồm chữ cái thường, và tổng \(n\) không vượt quá
\(20000\).

\subsection*{Output}

Với mỗi bộ dữ liệu, in \(n-1\) dòng. Dòng thứ \(i\) là số cặp cho cách cắt tại
\(i\).

\subsection*{Sample}

\begin{verbatim}
Input
3
4 0
jslj
6 1
abcazz
5 0
zzzzz

Output
1
1
1
5
9
10
8
5
4
8
8
4
\end{verbatim}

\section{Random Walk 2}

\paragraph{Time limit:} 1 second
\paragraph{Memory limit:} 256 megabytes

Penguin có một đồ thị có hướng đầy đủ gồm \(n\) đỉnh, có vòng lặp và không có
cạnh bội. Penguin đi ngẫu nhiên trên đồ thị:
\begin{enumerate}
  \item Bắt đầu tại đỉnh \(S\) ở thời điểm \(t=1\).
  \item Mỗi lần, nếu đang ở đỉnh \(i\), đi đến đỉnh \(j\) với xác suất
  \[
    P[i][j]=\frac{W[i][j]}{\sum_{k=1}^n W[i][k]}.
  \]
  \item Nếu tại thời điểm \(t\) và \(t+1\) đều ở đỉnh \(i\), Penguin sẽ ở lại
        đỉnh \(i\) mãi mãi.
\end{enumerate}
Cần tính xác suất \(A[i][j]\) để nếu bắt đầu tại \(i\) thì dừng lại tại \(j\).
Hãy tính \(A[i][j]\) modulo \(998244353\).

\subsection*{Input}

Dữ liệu vào gồm \(T\le 10\) bộ dữ liệu. Mỗi bộ dữ liệu gồm \(n\in[1,300]\), sau đó
là ma trận \(W\) kích thước \(n\times n\), với \(0\le W[i][j]\le 10^5\). Tổng
\(n\) không vượt quá \(2500\).

\subsection*{Output}

Với mỗi bộ dữ liệu, in \(n\) dòng, dòng thứ \(i\) gồm \(n\) số
\(A[i][j]\) modulo \(998244353\).

\subsection*{Sample}

\begin{verbatim}
Input
2
1
1
2
1 1
1 0

Output
665496236 332748118
332748118 665496236
1 0
1 0
\end{verbatim}

\section{Cute Tree}

\paragraph{Time limit:} 1 second
\paragraph{Memory limit:} 256 megabytes

Cho giả mã của hàm \(\operatorname{BuildTree}(A,id,L,R)\). Mỗi lần hàm được
gọi, biến \(\operatorname{tot}\) tăng lên \(1\), tạo một nút mới. Nếu
\(L=R\), hàm gán khóa của nút bằng \(A_L\). Nếu độ dài \([L,R]\) bằng \(1\),
hàm tách thành hai đoạn đơn \(L\) và \(R\). Ngược lại, đặt
\[
  B=L+\left\lceil\frac{\operatorname{len}(L,R)}{3}\right\rceil-1,\qquad
  C=\left\lfloor\frac{B+R}{2}\right\rfloor,
\]
rồi đệ quy trên \([L,B]\), \([B+1,C]\), \([C+1,R]\). Sau đó hàm duyệt
\(i\in[1,R]\) và cộng \(A_i\) vào khóa của nút.

Yêu cầu: tính số nút được tạo khi gọi
\[
  \operatorname{BuildTree}(A,\operatorname{root},1,n).
\]

\subsection*{Input}

Dữ liệu vào gồm \(T\) bộ dữ liệu (\(1\le T\le 5\)). Mỗi bộ dữ liệu gồm
\(n\) (\(1\le n\le 2\cdot 10^5\)) và mảng \(A_i\)
(\(1\le A_i\le 10^9\)).

\subsection*{Output}

Với mỗi bộ dữ liệu, in số nút được tạo.

\subsection*{Sample}

\begin{verbatim}
Input
1
7
4 3 5 2 6 7 1
2
4
2 2 5 3
10
21 10 5 89 12 3 42 13 55 76

Output
11
6
15
\end{verbatim}

\section{banzhuan}

\paragraph{Time limit:} 1 second
\paragraph{Memory limit:} 256 megabytes

Cho không gian ba chiều \([1,n]\times[1,n]\times[1,n]\). Cần đặt một số khối
lập phương \(1\times1\times1\) sao cho nhìn từ trên, từ trái và từ phía trước
đều thấy một hình vuông \(n\times n\). Mặt phẳng \(xOy\) là mặt đất, trục
\(z\) là độ cao.

Các khối phải tuân theo trọng lực: nếu bên dưới một khối là trống, khối đó sẽ
rơi xuống cho đến khi chạm đất (\(z=1\)) hoặc nằm trên một khối khác. Nếu đặt
một khối tại tọa độ nguyên \((x,y,z)\), chi phí của nó là
\[
  x\cdot y^2\cdot z.
\]
Trong mọi cách đặt thỏa điều kiện, hãy tính chi phí nhỏ nhất và lớn nhất.

\subsection*{Input}

Dữ liệu vào gồm \(T\le 15\) bộ dữ liệu. Mỗi bộ dữ liệu gồm một số
\[
  1\le n\le 10^{18}.
\]

\subsection*{Output}

Với mỗi bộ dữ liệu, in hai dòng: chi phí nhỏ nhất modulo \(10^9+7\), rồi chi
phí lớn nhất modulo \(10^9+7\).

\subsection*{Sample}

\begin{verbatim}
Input
1

Output
1
1
\end{verbatim}

\section{supermarket}

\paragraph{Time limit:} 2.5 seconds
\paragraph{Memory limit:} 256 megabytes

Có \(n\) loại hàng và \(m\) bản ghi mua hàng. Sau khi xóa thông tin lặp lại,
mỗi bản ghi được biểu diễn bằng một tập mặt hàng, mã hóa bởi số nguyên
\[
  A_i\in[0,2^n).
\]
Ví dụ \(n=4\), \(A_i=(0101)_2\) nghĩa là bản ghi chứa loại hàng thứ \(1\) và
thứ \(3\).

Định nghĩa \(P(T\mid S)\) là xác suất để một bản ghi ngẫu nhiên trong các bản
ghi chứa \(S\) cũng chứa \(T\). Nếu không có bản ghi nào chứa \(S\), thì
\(P(T\mid S)=0\) với mọi \(T\).

Cần tính
\[
  \sum_{S\in[0,2^n)}\sum_{T\in[0,2^n)} P(T\mid S)
\]
modulo \(998244353\).

\subsection*{Input}

Dữ liệu vào gồm \(T\le 15\) bộ dữ liệu. Mỗi bộ dữ liệu gồm
\[
  n\in[1,20],\qquad m\in[1,2\cdot10^5],
\]
sau đó là \(m\) số \(A_i\in[0,2^n)\).

\subsection*{Output}

Với mỗi bộ dữ liệu, in đáp án modulo \(998244353\).

\subsection*{Sample}

\begin{verbatim}
Input
3 5
7
6
5
3
5

Output
66549669
\end{verbatim}

\section{Array}

\paragraph{Time limit:} 15 seconds
\paragraph{Memory limit:} 512 megabytes

Cho mảng số nguyên \(a[1..n]\). Đếm số đoạn con \([L,R]\), \(R-L+1\ge1\),
sao cho tồn tại một giá trị mà số lần xuất hiện của nó lớn hơn tổng số lần
xuất hiện của tất cả giá trị khác trong \(a[L..R]\).

\subsection*{Input}

Dữ liệu vào gồm \(T\le 15\) bộ dữ liệu. Mỗi bộ dữ liệu gồm \(n\)
(\(1\le n\le 10^6\)) và mảng \(a_i\) với \(0\le a_i\le 10^6\).

\subsection*{Output}

Với mỗi bộ dữ liệu, in đáp án.

\subsection*{Sample}

\begin{verbatim}
Input
1
1
96705
1
10
3 7 3 3 3 7 3 3 7 7

Output
1
47
\end{verbatim}

\section{Guess Or Not 2}

\paragraph{Time limit:} 5 seconds
\paragraph{Memory limit:} 256 megabytes

Bob có một bộ sinh phân phối xác suất rời rạc từ một vector và tham số \(t\).
Bộ sinh lấy \(x_i\), biến đổi bằng logarit và hai lần số ngẫu nhiên đều trên
\((0,1)\), sau đó chuẩn hóa để thu được vector \(y\).

Cho vector đầu ra \(y[1..k]\), hãy tính mật độ xác suất của phân phối đầu ra
tại vector đó, ký hiệu
\[
  f_{\operatorname{Generate}(x,k,t)}(y),
\]
và in kết quả modulo \(998244353\).

\subsection*{Input}

Dữ liệu vào gồm \(T\) bộ dữ liệu. Mỗi bộ dữ liệu gồm \(k,t\), sau đó là \(k\) số
\(x_i\) và \(k\) số \(z_i\). Vector cần tính được cho bởi
\[
  y_i=\frac{z_i}{\sum_j z_j}.
\]
Tổng \(k\) không vượt quá \(10^6\), \(2\le k\le 10^6\), và
\[
  t,x_i,z_i\in[1,998244353).
\]
Đảm bảo \(\sum_i z_i\not\equiv0\pmod {998244353}\).

\subsection*{Output}

Với mỗi bộ dữ liệu, in một số nguyên là mật độ xác suất cần tính modulo
\(998244353\).

\subsection*{Sample}

\begin{verbatim}
Input
2
3 1
1 1 1
2 2 2
3 2
1 1 1
1 2 3

Output
2
596788047
\end{verbatim}

\section{jsljgame}

\paragraph{Time limit:} 1 second
\paragraph{Memory limit:} 256 megabytes

Có \(n\) đống đá, đống thứ \(i\) có \(a_i\) viên. Jslj và Penguin luân phiên
lấy đá, Jslj đi trước. Mỗi lượt, Jslj có thể lấy một số dương viên đá khác
\(x\) từ một đống bất kỳ; Penguin có thể lấy một số dương viên đá khác \(y\)
từ một đống bất kỳ. Người đầu tiên không thể đi sẽ thua. Hãy xác định người
thắng.

\subsection*{Input}

Dữ liệu vào gồm \(T\) bộ dữ liệu (\(1\le T\le 2000\)). Mỗi bộ dữ liệu gồm
\[
  n,x,y\qquad (1\le n\le 10^3,\ 1\le x,y\le 10^9),
\]
sau đó là \(a_i\) (\(1\le a_i\le 10^9\)).

\subsection*{Output}

Nếu Jslj thắng, in \texttt{Jslj}; ngược lại in \texttt{yygqPenguin}.

\subsection*{Sample}

\begin{verbatim}
Input
1
3 20 100
5 800000 10

Output
Jslj
\end{verbatim}

\section{Yet another matrix problem}

\paragraph{Time limit:} 4 seconds
\paragraph{Memory limit:} 256 megabytes

Có hai ma trận \(A\) và \(B\). Ma trận \(A_{n,r}\) có \(n\) hàng, \(r\) cột;
mỗi phần tử \(A[i][j]\) là số nguyên trong \([0,m]\). Ma trận \(B_{r,n}\) có
\(r\) hàng, \(n\) cột; mỗi phần tử \(B[i][j]\) cũng nằm trong \([0,m]\).
Định nghĩa \(f(x)\) là số cặp \((A_{n,r},B_{r,n})\) sao cho
\[
  C=A\times B,\qquad \sum_{i=1}^n\sum_{j=1}^n C[i][j]=x.
\]
Để đơn giản, đặt \(r=nm\). Cần tính
\[
  f(0),f(1),\ldots,f(m)\pmod {998244353}.
\]

\subsection*{Input}

Dữ liệu vào gồm \(T\le 15\) bộ dữ liệu. Mỗi bộ dữ liệu gồm hai số
\[
  n,m\in[1,10^5].
\]

\subsection*{Output}

Với mỗi bộ dữ liệu, in \(m+1\) dòng. Dòng thứ \(i\) in
\[
  f(i-1)\bmod 998244353.
\]

\subsection*{Sample}

\begin{verbatim}
Input
2 1

Output
49
56
\end{verbatim}

\section{penguin love tour}

\paragraph{Time limit:} 1.5 seconds
\paragraph{Memory limit:} 256 megabytes

Có \(n\) hòn đảo nối với nhau bằng \(n-1\) con đường vô hướng, đường thứ \(i\)
nối \(u_i,v_i\) và mất \(w_i\) thời gian để đi qua. Mỗi hòn đảo \(i\) có sức
mạnh \(p_i\). Với mỗi hòn đảo, ta được thực hiện tối đa một lần thao tác:
chọn một con đường kề với hòn đảo đó và đổi chi phí \(w\) của đường thành
\[
  \max(0,w-p_i).
\]
Penguin sẽ chọn tùy ý một hòn đảo và đi theo đường ngắn nhất đến một hòn đảo
khác. Vì Penguin thích đi xa, nó sẽ chọn đường dài nhất trong các lựa chọn có
thể. Hãy phân bổ sức mạnh của các hòn đảo vào các đường sao cho độ dài chuyến
đi dài nhất này nhỏ nhất.

\subsection*{Input}

Dữ liệu vào gồm \(T\le 1000\) bộ dữ liệu. Mỗi bộ dữ liệu gồm \(n\)
(\(1\le n\le 10^5\)), dãy \(p_i\) với \(0\le p_i\le 10^5\), và \(n-1\) cạnh
\((u_i,v_i,w_i)\), \(1\le w_i\le 10^5\). Tổng \(n\) không vượt quá
\(4\cdot10^5\).

\subsection*{Output}

Với mỗi bộ dữ liệu, in độ dài hành trình ngắn nhất có thể đạt được.

\subsection*{Sample}

\begin{verbatim}
Input
1
7
9 6 7 5 7 4 10
4 3 5
3 2 8
7 5 6
6 3 3
5 4 4
2 1 3
1
9
3 2 1 1 2 2 4 4 3
5 3 5
4 2 1
6 4 1
2 1 1
3 2 2
9 2 4
7 5 3
8 4 4

Output
0
3
\end{verbatim}

\end{document}
